\problemname{Duplantis}
\noindent

Duplantis has a problem. He can't decide if he should beat as many world records as possible, or set one world record that is as good 
as possible.

At the moment, his world record is $v$ centimeters. Duplantis has predicted how well he will perform in the remaining $N$ competitions
of his career. To do that, he calculated the positive integers $a_1, a_2, \dots, a_N$. The number $a_i$ means that Duplantis can jump between 
$0$ and $a_i$ centimeters in the $i$th competition. He can decide freely how high he should jump in a competition, as long as the height is 
an integer between $0$ and $a_i$ centimeters.

In order to beat the world record, Duplantis has to jump strictly higher than his previous world record. It is only allowed to beat at most 
one world record in the same competition. In pole vaulting, you always jump an integer number of centimeters.

Your task is to calculate $r$, the maximum number of additional world records that Duplantis can beat. You should also find the numbers 
$b_1, b_2, \dots, b_r$, where $b_i$ is the highest possible final world record given that Duplantis beats exactly $i$ more world records.

\section*{Input}

The first line contains two integers $N$ and $v$ ($1 \leq N \leq 3 \cdot 10^5$, $0 \leq v \leq 10^9$), where $N$ is the number of competitions
and $v$ is Duplantis' current world record. 

The second line contains $N$ integers $a_1, a_2, \dots, a_N$ ($0 \leq a_i \leq 10^9$). 


\section*{Output}

On the first line, print $r$, the maximum number of additional world records Duplantis can beat.

On the second line, print the $r$ numbers $b_1, b_2, \dots, b_r$, where $b_i$ is the maximum possible final world record if Duplantis
beats exactly $i$ more world records.

Note that $r$ can be zero, in that case you should leave the second line empty (it is also OK to not print this empty line).


\section*{Scoring}
Your solution will be tested on a set of test groups, each worth a number of points. Each test group contains
a set of test cases. To get the points for a test group you need to solve all test cases in the test group.

\noindent
\begin{tabular}{| l | l | p{12cm} |}
  \hline
  \textbf{Group} & \textbf{Points} & \textbf{Constraints} \\ \hline
  $1$    & $20$       & $v < a_1 < a_2 < \dots < a_N$ \\ \hline
  $2$    & $20$        & $a_1 > a_2 > \dots > a_N$ \\ \hline
  $3$    & $17$        & $N = 2$ \\ \hline
  $4$    & $43$       & No additional constraints. \\ \hline
\end{tabular}

\section*{Explanation of Samples}

In the first example, Duplantis can beat at most $3$ more world records. He can for example jump $635$ in the first competition, 
$640$ in the second, and finally $650$ in the fourth competition. In the other competitions he can do whatever he wants. 
On the other hand, if he only beats one more world record, then he can achieve a height of $700$ centimeters by jumping as high as possible
in the first competition.

In the second example it unfortunately seems like Duplantis has already peaked, and cannot beat any more world records.
